% Lesson 20 (AN-18): Retrieve into a document (the text goes in at the
% cursor; the document keeps its own name); Save under a new name.  Uses
% LETTERHD (Lesson 18) and PROCS (Lesson 9); saves PACKAGE.  Replayed by
% test-pubs-iword.
% GATES: test-pubs-iword
\Gates{test-pubs-iword}
\Lesson{20}{Combining Documents}
\LessonFeatures{Retrieve; Save; Filenames}
Retrieve onto a clear screen brings a document back.  Retrieve with a
document already on the screen puts the second document into the first, at
the cursor.  In this lesson you put the cover sheet procedure under the
letterhead of Lesson 18, and save the result as a new document.

\BeginReplay{l20}
\Task{Retrieving the Letterhead}
\begin{WPSteps}
\item Press \Press{Retrieve}{Shift-F10}, type \TypeName{letterhd} and press
\PressKey{Enter}.
\item Press \PressKey{Home}, \PressKey{Home}, \PressKey{Down} to go to the end,
under the salutation.
\end{WPSteps}

\Task{Retrieving the Procedure Into It}
\begin{WPSteps}
\item Press \Press{Retrieve}{Shift-F10}, type \TypeName{procs} and press
\PressKey{Enter}.
\end{WPSteps}
\CompareScreen[rows=17]{l20-1}
The procedure is now part of the letter, starting at the cursor, which is
still where it was.  The status line still shows \Scr{A:\wbs LETTERHD}: the
document on the screen keeps its name, and neither file on the disk has
changed.

\Task{Saving Under a New Name}
To keep the letterhead as it is, save the combined document under a new
name.
\begin{WPSteps}
\item Press \Press{Save}{F10}.  The status line offers \Scr{LETTERHD}.
\item Type \TypeName{package}.  The first character you type replaces the
name offered.  Press \PressKey{Enter}.
\end{WPSteps}
\CompareScreen[rows=17]{l20-2}

\Task{Clearing the Screen}
\begin{WPSteps}
\item Press \Press{Exit}{F7}, type \TypeName{n} and \TypeName{n}.
\end{WPSteps}
\EndReplay
In Lesson 21 you meet the messages InitechWord shows when something cannot
be done, and the way back from each.
